% Example for ai-math's LaTeX route, built from skills/ai-math/templates/article.tex.
% Build:  node skills/ai-math/scripts/latex-pdf.mjs examples/math/continuous-compounding.tex --preview
% Engine: Tectonic (XeTeX). Packages download on first use and are cached.
\documentclass[11pt]{article}

% ---------------------------------------------------------------- page
\usepackage[letterpaper, margin=1in, headheight=14pt]{geometry}
\usepackage{fontspec}                 % XeTeX: Unicode text, system or bundled fonts
\usepackage{unicode-math}             % OpenType math (Latin Modern Math by default)
\usepackage{microtype}
\usepackage[dvipsnames]{xcolor}
\usepackage{parskip}

% Seneca palette
\definecolor{ink}{HTML}{16191D}
\definecolor{steel}{HTML}{5E7A94}
\definecolor{signal}{HTML}{C9A000}   % darker than the screen yellow so it prints legibly
\color{ink}

% ---------------------------------------------------------------- math
\usepackage{amsmath, mathtools}
\usepackage{amsthm}
\newtheoremstyle{seneca}{10pt}{10pt}{}{}{\bfseries\sffamily\color{steel}}{.}{0.5em}{}
\theoremstyle{seneca}
\newtheorem{theorem}{Theorem}
\newtheorem{definition}{Definition}
\newtheorem{problem}{Problem}
\newcommand{\R}{\mathbb{R}}
\newcommand{\dd}{\mathop{}\!\mathrm{d}}

% ---------------------------------------------------------------- headings
\usepackage{titlesec}
\titleformat{\section}{\large\bfseries\sffamily}{\color{steel}\thesection}{0.6em}{}
\titleformat{\subsection}{\normalsize\bfseries\sffamily}{\color{steel}\thesubsection}{0.6em}{}
\usepackage{fancyhdr}
\pagestyle{fancy}\fancyhf{}
\renewcommand{\headrulewidth}{0pt}
\fancyfoot[C]{\small\color{steel}\thepage}

\usepackage[hidelinks]{hyperref}

\newcommand{\kicker}[1]{{\sffamily\footnotesize\bfseries\color{steel}\MakeUppercase{#1}}}

\begin{document}

\kicker{Money math \quad·\quad derivation}\par
{\sffamily\bfseries\fontsize{24}{28}\selectfont Where $e$ comes from in compound interest\par}
\vspace{2pt}
{\color{steel}A companion to the one-page compound interest sheet}

\section{The question}
A balance $P$ earns an annual rate $r$, compounded $n$ times a year. After $t$ years it is
\begin{equation}\label{eq:discrete}
  A_n(t) = P\left(1 + \frac{r}{n}\right)^{nt}.
\end{equation}
Compounding more often always helps, but by less each time. Does $A_n(t)$ grow without bound as $n \to \infty$,
or does it settle?

\begin{definition}
Euler's number is the limit $\displaystyle e = \lim_{m\to\infty}\left(1 + \frac{1}{m}\right)^{m} \approx 2.71828$.
\end{definition}

\section{The limit}
\begin{theorem}\label{thm:continuous}
For every $r > 0$ and $t \ge 0$, \quad $\displaystyle \lim_{n\to\infty} A_n(t) = P e^{rt}$.
\end{theorem}

\begin{proof}
Substitute $m = n/r$, so $r/n = 1/m$ and $nt = m\,rt$. Then
\begin{align*}
  A_n(t) &= P\left(1 + \frac{1}{m}\right)^{m\,rt}
          = P\left[\left(1 + \frac{1}{m}\right)^{m}\right]^{rt}.
\end{align*}
As $n \to \infty$, $m \to \infty$ too. The inner bracket tends to $e$, and $x \mapsto x^{rt}$ is continuous,
so the whole expression tends to $P e^{rt}$.
\end{proof}

\section{How fast the gap closes}
Taking logarithms and expanding $\ln(1+x) = x - \tfrac{x^2}{2} + O(x^3)$ with $x = r/n$:
\begin{equation}
  \ln A_n(t) = \ln P + nt\left(\frac{r}{n} - \frac{r^2}{2n^2} + O(n^{-3})\right)
             = \ln P + rt - \frac{r^2 t}{2n} + O(n^{-2}).
\end{equation}
So the shortfall from continuous compounding is roughly a factor $e^{-r^2 t / 2n}$: halving it means doubling $n$.

\begin{problem}
At $r = 5\%$ for $t = 30$ years, how much more does continuous compounding pay than monthly ($n = 12$) on \$10{,}000?
\end{problem}
\emph{Answer.} $10{,}000\,e^{1.5} - 10{,}000\left(1 + \tfrac{0.05}{12}\right)^{360} \approx 44{,}816.89 - 44{,}677.44 \approx \$139$,
in line with $e^{-r^2t/2n} = e^{-0.003125} \approx 0.9969$ from the estimate above.

\end{document}
